%! Setting Up --- Learning from Data — Unit 8, Lesson 8.0 — Slides (Beamer)
\documentclass[aspectratio=169, 11pt]{beamer}
\usepackage{linalg-beamer}   % provides \burgundyheader{} and \sectionlabel[color]{}
\newcommand{\vv}{\mathbf{v}}
\newcommand{\bb}{\mathbf{b}}
\newcommand{\relu}{\mathrm{ReLU}}

\begin{document}

% TITLE SLIDE (CourseName is not defined in beamer, so the name is literal)
{
\setbeamercolor{background canvas}{bg=burgundy}
\begin{frame}[plain]
  \vspace{2.8cm}\hspace{0.55cm}%
  \begin{minipage}{0.9\textwidth}
    {\color{goldacc}\bfseries\small LESSON 8.0}\\[0.5cm]
    {\color{white}\bfseries\LARGE Setting Up --- Learning from Data}\\[1.0cm]
    {\color{blushmid}\small Linear Algebra~$\cdot$~Unit 8: Learning from Data}
  \end{minipage}
\end{frame}
}

% HOOK
\begin{frame}[t, plain]
  \burgundyheader{The linear algebra you know is the engine of AI}
  \sectionlabel{Hook}
  \begin{columns}[T]
    \begin{column}{0.5\textwidth}
      A \textbf{neural network} is just a \emph{function} that turns an input into a prediction. It is
      built from things you already have:\\[6pt]
      a \textbf{linear step} --- multiply by a matrix, add a bias (Unit~1) --- and \emph{one} new
      nonlinear step, the \textbf{ReLU} $=\max(0,x)$.
    \end{column}
    \begin{column}{0.5\textwidth}
      \centering
      \begin{tikzpicture}[scale=0.5]
        \draw[->, charcoal, thin] (-2.2,0)--(2.5,0) node[right]{\scriptsize $x$};
        \draw[->, charcoal, thin] (0,-0.5)--(0,2.6) node[above]{\scriptsize $\relu(x)$};
        \draw[ultra thick, burgundy] (-2,0)--(0,0)--(2,2);
        \fill[burgundy] (0,0) circle (1.6pt);
        \node[charcoal] at (1.0,-0.9) {\scriptsize one fold at $0$};
      \end{tikzpicture}
    \end{column}
  \end{columns}
  \vspace{2pt}
  \begin{block}{The question of Unit 8}
    How do a matrix and one simple bend combine into a machine that \textbf{learns from data}?
  \end{block}
\end{frame}

% RELU
\begin{frame}[t, plain]
  \burgundyheader{The one new ingredient: ReLU $=\max(0,x)$}
  \sectionlabel[goldacc]{Definition}
  \begin{itemize}
    \setlength{\itemsep}{6pt}
    \item The \textbf{ReLU} keeps positives and replaces negatives with $0$:
          \[ \relu(3)=3,\quad \relu(-2)=0,\quad \relu(0)=0. \]
    \item As a graph it is one \textbf{fold} --- flat, then a line. Add and shift a few and their folds
          combine into a \textbf{piecewise-linear} function that can \emph{bend} to fit data.
    \item Example: $F(x)=\relu(x)-\relu(x-2)$ rises from $0$ to $2$ over $0\le x\le2$, then stays flat ---
          two folds from two ReLUs.
  \end{itemize}
\end{frame}

% ONE LAYER
\begin{frame}[t, plain]
  \burgundyheader{A learning function is a chain of matrix steps}
  \sectionlabel{Skill}
  \begin{block}{One layer}
    \[ \vv \;\longmapsto\; \relu(A\vv+\bb): \quad\text{linear step (Unit 1), then bend with ReLU.} \]
  \end{block}
  \begin{columns}[T]
    \begin{column}{0.5\textwidth}
      With $A=\begin{bsmallmatrix} 1 & -1\\ 1 & 1 \end{bsmallmatrix}$, $\bb=\mathbf{0}$,
      $\vv=\begin{bsmallmatrix} 1\\ 2 \end{bsmallmatrix}$:
      \[ A\vv=\begin{bsmallmatrix} -1\\ 3 \end{bsmallmatrix},\qquad
         \relu(A\vv)=\begin{bsmallmatrix} 0\\ 3 \end{bsmallmatrix}. \]
      The ReLU zeroed the negative entry.
    \end{column}
    \begin{column}{0.5\textwidth}
      \textbf{Matrix multiplication does the heavy lifting} in every layer. A network is mostly $A\vv$,
      repeated, with a ReLU in between. The bias $\bb$ just shifts where each fold sits.
    \end{column}
  \end{columns}
\end{frame}

% HOW IT LEARNS + ROAD AHEAD
\begin{frame}[t, plain]
  \burgundyheader{How it learns --- and the road ahead}
  \sectionlabel[goldacc]{Payoff}
  \begin{itemize}
    \setlength{\itemsep}{6pt}
    \item \textbf{Loss:} the Unit~4 sum of squared errors between predictions and targets. Predict
          $3,5,4$ vs.\ targets $2,5,6$: loss $=1+0+4=5$.
    \item \textbf{Learning $=$ gradient descent:} nudge the weights to roll the loss \emph{downhill}
          (Lesson~8.3). The data's mean/variance/covariance describe its shape (Lesson~8.4).
    \item \textbf{Why it matters:} a neural network is matrix multiplication (Unit~1) $+$ one bend (ReLU),
          trained by least squares (Unit~4). \emph{No new machinery.}
  \end{itemize}
  \vspace{2pt}
  \begin{block}{The road ahead}
    \textbf{8.1} piecewise-linear learning functions \;$\bullet$\; \textbf{8.2} convolutional neural nets
    \;$\bullet$\; \textbf{8.3} minimizing loss by gradient descent \;$\bullet$\; \textbf{8.4} mean,
    variance, covariance.
  \end{block}
\end{frame}

\end{document}
